\documentclass[tikz,border=15pt]{standalone}
\usepackage[utf8]{inputenc}
\usepackage[spanish]{babel}
\usepackage{helvet}
\renewcommand{\familydefault}{\sfdefault}
\usetikzlibrary{shapes.geometric,arrows.meta,positioning,calc,fit,shadows}

\definecolor{azul}{HTML}{173B59}
\definecolor{verde}{HTML}{176B66}
\definecolor{suave}{HTML}{EDF3F7}
\definecolor{gris}{HTML}{46515A}
\definecolor{linea}{HTML}{CAD6DF}

\begin{document}
\begin{tikzpicture}[
  x=1.0cm, y=1.0cm,
  raiz/.style={
    draw=azul,fill=azul,text=white,rounded corners=5pt,
    text width=3.8cm,align=center,inner sep=10pt,
    font=\small\bfseries,drop shadow
  },
  rama/.style={
    draw=verde,fill=suave,rounded corners=4pt,
    text width=3.2cm,align=left,inner sep=7pt,
    font=\small\bfseries
  },
  detalle/.style={
    draw=linea,fill=white,rounded corners=4pt,
    text width=14.0cm,align=left,inner sep=7pt,
    font=\footnotesize
  },
  enlace/.style={draw=verde,line width=1.1pt}
]
\node[raiz] (r) at (0,0) {ORDEN JUR\'IDICO\\NACIONAL EN\\DERECHOS HUMANOS\\[5pt]
\footnotesize Sistema articulado de fuentes, principios y garant\'ias};

\node[rama] (a) at (4.5, 4.5) {1. Concepto y fines};
\node[detalle] (ad) at (13.8, 4.5) {
\textbf{Concepto:} Conjunto arm\'onico de normas y competencias constitucionales y convencionales.\\
\textbf{Finalidad y eficacia:} Tutelar la dignidad intr\'inseca y regular el poder p\'ublico como par\'ametro vinculante de validez.
};

\node[rama] (b) at (4.5, 2.7) {2. Fundamento constitucional};
\node[detalle] (bd) at (13.8, 2.7) {
\textbf{Bloque constitucional (Art. 1.º):} Cl\'ausula de apertura internacional y obligaciones para todas las autoridades.\\
\textbf{Supremac\'ia y garant\'ias (Arts. 102 Ap. B, 103, 107 y 133):} Juicio de amparo y sistema ombudsperson no jurisdiccional.
};

\node[rama] (c) at (4.5, 0.9) {3. Principios rectores};
\node[detalle] (cd) at (13.8, 0.9) {
\textbf{Universalidad e interdependencia:} Titularidad incondicional de toda persona e impacto rec\'iproco entre derechos.\\
\textbf{Indivisibilidad y progresividad:} Integridad del haz de prerrogativas, avance continuo y prohibici\'on de retroceso.
};

\node[rama] (d) at (4.5, -0.9) {4. Fuentes internacionales};
\node[detalle] (dd) at (13.8, -0.9) {
\textbf{Sistemas internacional y regional:} Convenci\'on Americana (Corte IDH), Pactos ONU 1966 (PIDCP y PIDESC).\\
\textbf{Tratados tem\'aticos y DUDH:} Protecci\'on frente a tortura, discriminaci\'on y directrices hermen\'euticas universales.
};

\node[rama] (e) at (4.5, -2.7) {5. Legislaci\'on de desarrollo};
\node[detalle] (ed) at (13.8, -2.7) {
\textbf{Leyes reglamentarias:} Ley de Amparo, Ley de la CNDH y Ley General de V\'ictimas (reparaci\'on integral y asesor\'ia).\\
\textbf{Leyes transversales:} Normas de no discriminaci\'on, igualdad de g\'enero y acceso a la informaci\'on p\'ublica.
};

\node[rama] (f) at (4.5, -4.5) {6. Instancias de garant\'ia};
\node[detalle] (fd) at (13.8, -4.5) {
\textbf{Jurisdiccional y ombudsperson:} Poder Judicial de la Federaci\'on (SCJN, TCC, Juzgados de Distrito) y CNDH/CEDH.\\
\textbf{Control difuso y sede subsidiaria:} Deber de todas las autoridades locales y recurso ante la Comisi\'on y Corte IDH.
};

\draw[enlace] (r.east) -- (2.6,0);
\draw[enlace] (2.6, -4.5) -- (2.6, 4.5);
\foreach \n in {a,b,c,d,e,f}{
  \draw[enlace] (2.6,0 |- \n.west) -- (\n.west);
}
\foreach \n/\m in {a/ad,b/bd,c/cd,d/dd,e/ed,f/fd}{
  \draw[enlace] (\n.east) -- (\m.west);
}
\end{tikzpicture}
\end{document}
